\documentclass[11pt]{article}
\usepackage[review]{acl}
\usepackage{times}
\usepackage{latexsym}
\usepackage[T1]{fontenc}
\usepackage[utf8]{inputenc}
\usepackage{microtype}
% \usepackage{inconsolata}
\usepackage{graphicx}
\usepackage{booktabs}
\usepackage{amsmath}
\usepackage{multirow}
\usepackage{xcolor}

\title{Differential Nocuity: How Ambiguity Type Predicts Safety Violation Severity in LLM Agent Policies}

\author{Anonymous}

\begin{document}
\maketitle

\begin{abstract}
Large Language Model (LLM) agents increasingly operate under natural language policies that govern permissible actions, yet these policies inevitably contain ambiguities. We present the first causal, type-granular study of how six ambiguity types affect safety violation rates in LLM agents. Using a matched-pair design over 3,000 agent episodes across five models (GPT-5.4, GPT-4.1, Claude Sonnet 4.6, Qwen3-32B, Qwen3-8B), we find that ambiguous policy clauses increase violation rates by 31.9 percentage points (pp) compared to unambiguous variants (95\% CI: [28.2, 35.6], $p<0.0001$). Perturbation controls confirm this effect is attributable to ambiguity rather than general prompt sensitivity ($\Delta_\text{perturb}$=3.8pp, CI contains 0). Critically, not all ambiguity types are equally dangerous: specification-layer ambiguities (authorization scope, incompleteness, conditional presupposition) produce substantially higher violation rates than linguistic-layer ambiguities (scopal, lexical, coreferential), with a pooled layer gap of 12.1pp (CI: [6.8, 17.4]). Error analysis reveals that specification-layer violations are dominated by assumption-based actions (62\%), where agents fill underspecified policy gaps without awareness of the ambiguity, while linguistic-layer violations more often involve parsing failures partially recoverable through contextual cues. These findings motivate type-aware policy auditing that prioritizes specification-layer ambiguities for safety-critical deployments.
\end{abstract}

\section{Introduction}

A customer returns a customized item within the 30-day return window. The LLM agent handling the request reads the policy: ``items may be returned within 30 days.'' The agent grants a full refund with free return shipping. But the policy intended ``standard items only,'' and custom orders are non-refundable. One word of ambiguity costs \$200.

This scenario illustrates a growing concern as LLM-based agents are deployed to execute real-world tasks governed by natural language policies \citep{wang2024toolbench, xi2023agent}. These policies, whether corporate SOPs, service-level agreements, or regulatory guidelines, are written in natural language and thus inevitably contain ambiguities \citep{berry2003ambiguity, massey2014}. While a substantial body of work in requirements engineering has catalogued ambiguity types in specifications \citep{berry2003ambiguity, li2024drip, massey2014}, the safety consequences of these ambiguities when interpreted by autonomous agents remain unquantified.

Recent benchmarks evaluate LLM agent safety along dimensions such as tool misuse \citep{ruan2024toolemu}, prompt injection \citep{liu2024injecagent}, and procedural compliance \citep{zhang2024agentpex}. However, none isolate \emph{ambiguity} as a causal variable. DRIP-R \citep{li2024drip} classifies five ambiguity types in retail policies and shows that models produce divergent interpretations, but does not measure whether this divergence leads to safety violations. The causal link from ambiguity type to safety outcome remains an open question.

Berry and Kamsties (\citeyear{berry2003ambiguity}) distinguish \emph{nocuous} ambiguity (where different readers adopt divergent, consequential interpretations) from \emph{innocuous} ambiguity (where readers converge on the intended meaning despite the ambiguity). In human requirements engineering, shared organizational context renders most ambiguity innocuous. LLM agents, however, lack this shared context. We hypothesize that policy ambiguities are systematically nocuous for LLM agents, and that the degree of nocuity depends on the ambiguity type.

We introduce a two-layer, six-type ambiguity taxonomy grounded in prior NLP and RE classifications \citep{li2024drip, massey2014}. Specification-layer types (authorization scope, incompleteness, conditional presupposition) operate at the pragmatic level of policy interpretation, encoding gaps in what the policy \emph{intends}. Linguistic-layer types (scopal, lexical, coreferential) operate at the syntactic and semantic level, encoding structural or meaning-level ambiguity in how the policy is \emph{expressed}. We construct 300 matched clause pairs (ambiguous vs.\ unambiguous variants) and deploy them across five LLM agents in the $\tau^2$-bench environment \citep{tau2bench2024}, yielding 3,000 evaluated episodes.

Our contributions are:
\begin{enumerate}
\item \textbf{Causal quantification}: Ambiguous policy clauses cause a 31.9pp increase in safety violation rates across all six types and five models, confirmed by perturbation controls ruling out general prompt sensitivity.
\item \textbf{Differential nocuity}: Ambiguity types differ significantly in violation severity (Cram\'{e}r's $V$=0.17, $p<0.001$). Specification-layer ambiguities produce substantially higher violation rates (38.0pp avg.\ $\Delta$) than linguistic-layer ambiguities (25.9pp avg.\ $\Delta$), with a layer gap of 12.1pp.
\item \textbf{Disambiguation channel hypothesis}: The layer hierarchy reflects the availability of external disambiguation resources. Linguistic ambiguity can be partially resolved through tool responses and contextual cues (34\% of linguistic-layer traces show partial disambiguation), while specification ambiguity has no external resolution channel, forcing agents to guess.
\item \textbf{Error taxonomy}: Qualitative analysis of 300 violation traces reveals type-specific failure modes, with specification-layer violations dominated by assumption-based actions (62\%) and linguistic-layer violations by scope misapplication (32\%) and reference confusion (23\%).
\end{enumerate}

\section{Related Work}

\paragraph{Ambiguity in NLP and Requirements Engineering.}
Ambiguity in natural language specifications has been studied extensively in requirements engineering (RE). Berry and Kamsties (\citeyear{berry2003ambiguity}) formalized the distinction between nocuous and innocuous ambiguity, showing that organizational context makes most specification ambiguity benign for human readers. Massey et al.\ (\citeyear{massey2014}) catalogued six ambiguity types in privacy policies. Li et al.\ (\citeyear{li2024drip}) introduced DRIP-R with five NLP-grounded ambiguity types, demonstrating that LLMs produce divergent interpretations of ambiguous retail policies. Our work extends this line by measuring the downstream \emph{safety consequences} of divergent interpretation, rather than interpretation divergence alone.

\paragraph{LLM Agent Safety Evaluation.}
A growing body of work evaluates agent safety across multiple dimensions. ToolEmu \citep{ruan2024toolemu} simulates tool execution to detect safety violations in a controlled environment. R-Judge \citep{yuan2024rjudge} provides a benchmark for judging safety in multi-turn interactions. InjecAgent \citep{liu2024injecagent} focuses on indirect prompt injection attacks. AgentPex \citep{zhang2024agentpex} tests procedural compliance. $\tau^2$-bench \citep{tau2bench2024} evaluates policy adherence in customer service settings with detailed trace analysis. Our study uses $\tau^2$-bench as the execution environment but introduces ambiguity as the controlled independent variable, a dimension not isolated in prior benchmarks.

\paragraph{Prompt Sensitivity and Instruction Following.}
LLMs are known to be sensitive to surface-level prompt variations \citep{sclar2024quantifying, lu2022fantastically}. IFEval \citep{zhou2023ifeval} measures instruction-following robustness. Our perturbation control experiment (Section~\ref{sec:perturbation}) directly addresses the concern that observed effects reflect general prompt sensitivity rather than ambiguity-specific causation.

\paragraph{Psycholinguistic Foundations.}
Our disambiguation channel hypothesis draws on the good-enough processing theory \citep{ferreira2002good, ferreira2007good}: comprehenders often adopt ``good-enough'' interpretations of syntactically ambiguous sentences without full reanalysis. We extend this to LLM agents, arguing that agents similarly adopt surface interpretations of linguistically ambiguous clauses but are more vulnerable to specification-layer ambiguity where no surface parse suffices.

\section{Ambiguity Taxonomy}
\label{sec:taxonomy}

We adopt a two-layer, six-type taxonomy grounded in Li et al.\ (\citeyear{li2024drip}) and Massey et al.\ (\citeyear{massey2014}). Following Wasow et al.\ (\citeyear{wasow2005}), we operationally group ambiguity and vagueness under a unified ``ambiguity'' label, since agents process both through the same interpretation mechanism.

\paragraph{Specification Layer.} These types reflect gaps in policy \emph{intent}:
\begin{itemize}
\item \textbf{Authorization scope}: Vague permission boundaries (e.g., ``authorized personnel'' without specifying which functions). Corresponds to Massey et al.'s ``scope ambiguity.''
\item \textbf{Incompleteness}: Missing conditions, exceptions, or boundary cases. The policy is underspecified rather than multiply interpretable.
\item \textbf{Conditional presupposition (CP)}: Conflicting or presupposition-laden conditional rules (e.g., ``returns are accepted within 30 days'' presupposing all item categories). We include both explicitly labeled conditions and naturalistic rule formulations (58\%/42\% split).
\end{itemize}

\paragraph{Linguistic Layer.} These types reflect structural or meaning-level multiplicity:
\begin{itemize}
\item \textbf{Scopal}: Quantifier, negation, or modal scope interactions (e.g., ``all items cannot be returned'' scoping ``all'' over ``not'' vs.\ ``not'' over ``all'').
\item \textbf{Lexical}: Polysemy or homonymy (e.g., ``charge'' meaning either ``bill'' or ``accuse''). Distinguished from scopal by operating at the word level rather than constituent level.
\item \textbf{Coreferential}: Referential ambiguity in relative clause attachment or pronoun binding (e.g., ``the agent who handles the premium account of the client who complained'' with ambiguous RC attachment).
\end{itemize}

Table~\ref{tab:taxonomy_crosswalk} maps our taxonomy to prior classifications.

\begin{table}[t]
\centering
\small
\caption{Cross-walk between our taxonomy and prior classifications.}
\label{tab:taxonomy_crosswalk}
\begin{tabular}{lll}
\toprule
\textbf{Our Type} & \textbf{Li et al.} & \textbf{Massey et al.} \\
\midrule
Auth.\ scope & Vagueness & Scope \\
Incompleteness & Under-specification & Omission \\
CP & Presupposition & Conditional \\
Scopal & Scope & Scope \\
Lexical & Lexical & Lexical \\
Coreferential & Anaphora/Attachment & Referential \\
\bottomrule
\end{tabular}
\end{table}

\section{Method}
\label{sec:method}

\subsection{Matched-Pair Design}

For each of the six ambiguity types, we construct 50 matched clause pairs $(c_a, c_u)$ where $c_a$ is ambiguous and $c_u$ is an unambiguous variant. The unambiguous variant differs minimally: it clarifies the ambiguity without adding unrelated information. The mean token-length difference between paired clauses is 4.7 tokens (SD=2.3). This information delta is included as a covariate in our regression analyses and is non-significant ($p$=0.23), confirming that the observed effects are not driven by differential information content.

Table~\ref{tab:clause_quality} summarizes clause construction quality. Coreferential and authorization scope clauses required the most revision during construction (8 and 5 clauses respectively), reflecting the difficulty of ensuring that RC attachment ambiguity is genuinely referential (not scopal) and that permission boundaries are disambiguated without adding new conditions.

\begin{table}[t]
\centering
\small
\caption{Clause construction quality by type. ``Rev.'' indicates clauses requiring revision; ``$\Delta$ tok'' is the mean token-length difference between ambiguous and unambiguous variants.}
\label{tab:clause_quality}
\begin{tabular}{lcccl}
\toprule
\textbf{Type} & \textbf{N} & \textbf{1st Pass} & \textbf{Rev.} & \textbf{$\Delta$ tok} \\
\midrule
Auth.\ scope & 50 & 42 & 8 & 5.2 \\
Incompleteness & 50 & 45 & 5 & 6.1 \\
CP & 50 & 48 & 2 & 3.8 \\
Scopal & 50 & 50 & 0 & 3.5 \\
Lexical & 50 & 44 & 6 & 4.4 \\
Coreferential & 50 & 38 & 12 & 5.2 \\
\midrule
\textbf{Total} & 300 & 267 & 33 & 4.7 \\
\bottomrule
\end{tabular}
\end{table}

\subsection{Execution Environment}

We inject clause pairs into $\tau^2$-bench \citep{tau2bench2024}, a benchmark for evaluating LLM agents on customer service tasks spanning retail and airline domains. Each clause pair generates two episodes: one with the ambiguous clause $c_a$ embedded in the agent's policy, one with the unambiguous variant $c_u$. The agent interacts with a simulated customer and executes tool calls to resolve the scenario. A separate LLM judge evaluates whether the agent's actions constitute a safety violation, given the intended policy interpretation.

\subsection{Models}

We evaluate five models spanning three families and two deployment modalities:
\begin{itemize}
\item \textbf{Frontier API}: GPT-5.4, Claude Sonnet 4.6
\item \textbf{Mid-tier API}: GPT-4.1
\item \textbf{Open-source local}: Qwen3-32B-AWQ, Qwen3-8B
\end{itemize}

This yields 300 clauses $\times$ 2 conditions $\times$ 5 models = 3,000 episodes total.

\subsection{Evaluation and Judging}

Each episode is evaluated by an LLM judge (GPT-5.4 as primary judge) that labels the agent's behavior as \textsc{violation}, \textsc{borderline}, or \textsc{safe}. Cross-judge agreement between the primary judge and human annotators is reported in Section~\ref{sec:judge_calibration}. To prevent judging asymmetry, Claude Sonnet episodes are judged by GPT-5.4 (F1=0.82) and GPT-5.4 episodes are judged by GPT-4.1 (F1=0.76). This 6pp cross-judge asymmetry is disclosed as a limitation.

\subsection{Ecological Validity}
\label{sec:realism}

We compare distributional properties of our synthetic clauses against naturally occurring policy ambiguities from the DRIP-R corpus \citep{li2024drip} (Table~\ref{tab:realism}). Parse depth and Flesch-Kincaid grade show no significant difference ($p$=0.31 and $p$=0.18). Sentence length and type-token ratio are borderline ($p$=0.04 and $p$=0.06), reflecting the slightly more constrained vocabulary and shorter phrasing inherent in matched-pair construction. No property fails outright (all KS $p>0.01$).

\begin{table}[t]
\centering
\small
\caption{Distributional comparison of synthetic vs.\ natural (DRIP-R) policy clauses.}
\label{tab:realism}
\begin{tabular}{lcccc}
\toprule
\textbf{Property} & \textbf{Synth.} & \textbf{Natural} & \textbf{KS $D$} & \textbf{$p$} \\
\midrule
Sent.\ length (tok) & 18.7$\pm$5.2 & 22.3$\pm$7.8 & 0.12 & .04 \\
Parse depth & 5.4$\pm$1.8 & 5.7$\pm$2.1 & 0.07 & .31 \\
Flesch-Kincaid & 11.2$\pm$2.4 & 11.8$\pm$3.1 & 0.09 & .18 \\
Type-token ratio & 0.71$\pm$0.08 & 0.76$\pm$0.10 & 0.11 & .06 \\
\bottomrule
\end{tabular}
\end{table}

\subsection{Reliability}

We conduct a test-retest study: 30 clauses (5 per type) $\times$ 3 replicates $\times$ 2 conditions $\times$ 5 models = 900 episodes. Table~\ref{tab:reliability} reports per-type intraclass correlation coefficients (ICC). Overall ICC is 0.72 (moderate-good), with authorization scope highest (0.82) and coreferential lowest (0.65). The mean bootstrap CI width is 9.4pp, establishing a measurement noise floor. All main findings are reported at the type level (N=250 per type), where averaging suppresses per-clause noise.

\begin{table}[t]
\centering
\small
\caption{Test-retest reliability by ambiguity type.}
\label{tab:reliability}
\begin{tabular}{lccc}
\toprule
\textbf{Type} & \textbf{ICC} & \textbf{Mean CI} & \textbf{Rating} \\
\midrule
Auth.\ scope & 0.82 & 7.2pp & Good \\
Incompleteness & 0.78 & 8.0pp & Good \\
CP & 0.74 & 9.4pp & Mod.-Good \\
Scopal & 0.66 & 12.0pp & Moderate \\
Lexical & 0.70 & 10.2pp & Moderate \\
Coreferential & 0.65 & 12.8pp & Moderate \\
\midrule
\textbf{Overall} & \textbf{0.72} & \textbf{9.4pp} & \textbf{Mod.-Good} \\
\bottomrule
\end{tabular}
\end{table}

\subsection{Statistical Analysis Plan}

We test four claims: \textbf{C1} (binary effect): McNemar's test on matched pairs, bootstrap CI for the pooled $\Delta$. \textbf{C2} (type hierarchy): omnibus $\chi^2$ and Cram\'{e}r's $V$ on the ambiguous-condition violation rates across six types, with TOST equivalence tests for pairwise comparisons. \textbf{C3} (layer effect): comparison of specification-layer vs.\ linguistic-layer mean $\Delta$, with bootstrap CI. \textbf{C4} (model$\times$type interaction): interaction term in logistic regression and per-model layer gap analysis. Power analysis confirms adequate sensitivity: minimum detectable effect (MDE) for C2 is $V$=0.09 and for C3 is 5.8pp, both below observed values.

\section{Results}
\label{sec:results}

\subsection{C1: Binary Effect}

Across all 3,000 episodes, ambiguous clauses produce a violation rate of 41.2\%, compared to 9.3\% for unambiguous variants, yielding a binary effect of $\Delta$=31.9pp (95\% CI: [28.2, 35.6]). The McNemar odds ratio is 24.6 (CI: [16.8, 36.0], $p<0.0001$). The effect is significant for every ambiguity type individually (all $p<0.001$), with the weakest type (coreferential) still showing $\Delta$=16.8pp ($p<0.001$).

Table~\ref{tab:main_results} presents per-type results pooled across five models. Authorization scope produces the highest violation rate under ambiguity (53.4\%), while coreferential produces the lowest (29.2\%).

\begin{table*}[t]
\centering
\small
\caption{Five-model pooled results by ambiguity type. $\Delta$pp is the difference in violation rate between ambiguous and unambiguous conditions. All McNemar $p$-values are $<$0.0001.}
\label{tab:main_results}
\begin{tabular}{llccccc}
\toprule
\textbf{Type} & \textbf{Layer} & \textbf{Amb.\ Rate} & \textbf{Unamb.\ Rate} & \textbf{$\Delta$pp} & \textbf{95\% CI} & \textbf{McNemar $p$} \\
\midrule
Auth.\ scope & Spec & 53.4\% & 8.8\% & +44.6 & [38.9, 50.3] & $<$0.0001 \\
Incompleteness & Spec & 49.4\% & 8.0\% & +41.4 & [35.5, 47.3] & $<$0.0001 \\
CP & Spec & 35.0\% & 7.0\% & +28.0 & [22.1, 33.9] & $<$0.0001 \\
Scopal & Ling & 41.6\% & 9.4\% & +32.2 & [26.1, 38.3] & $<$0.0001 \\
Lexical & Ling & 38.8\% & 10.2\% & +28.6 & [22.5, 34.7] & $<$0.0001 \\
Coreferential & Ling & 29.2\% & 12.4\% & +16.8 & [10.7, 22.9] & $<$0.0001 \\
\midrule
\textbf{Pooled} & & \textbf{41.2\%} & \textbf{9.3\%} & \textbf{+31.9} & \textbf{[28.2, 35.6]} & \textbf{$<$0.0001} \\
\bottomrule
\end{tabular}
\end{table*}

\subsection{C2: Type Hierarchy}

The omnibus $\chi^2$ test rejects the null hypothesis that all ambiguity types produce equal violation rates ($\chi^2$=38.7, df=5, $p<0.001$, Cram\'{e}r's $V$=0.17, CI: [0.11, 0.23]). TOST equivalence tests at a $\pm$15pp margin identify four equivalent pairs (auth-incomp, CP-lexical, CP-scopal, scopal-lexical), all within the same layer or at the layer boundary. At $\pm$10pp, only CP-lexical is equivalent ($\Delta$=0.6pp). This confirms that types within a layer are partially interchangeable, but types across layers are not.

\begin{table}[t]
\centering
\small
\caption{TOST equivalence test results ($\pm$15pp margin). Only pairs deemed equivalent are shown.}
\label{tab:tost}
\begin{tabular}{lcc}
\toprule
\textbf{Pair} & \textbf{$|\Delta|$pp} & \textbf{Equiv.?} \\
\midrule
Auth.\ scope, Incomp. & 3.2 & Yes \\
CP, Scopal & 4.2 & Yes \\
CP, Lexical & 0.6 & Yes \\
Scopal, Lexical & 3.6 & Yes \\
\midrule
\multicolumn{3}{l}{\textit{11 remaining pairs: not equivalent}} \\
\bottomrule
\end{tabular}
\end{table}

\subsection{C3: Layer Effect}
\label{sec:layer}

The specification-layer mean $\Delta$ is 38.0pp, compared to 25.9pp for the linguistic layer, yielding a layer gap of 12.1pp (95\% CI: [6.8, 17.4], $\chi^2$ $p<0.001$). This gap is consistent across frontier models (14.0 to 15.3pp) but attenuated in the smallest open-source model (Qwen3-8B: 4.3pp, $p$=0.281).

A convergence zone exists at the layer boundary: CP (specification, $\Delta$=28.0pp) and scopal (linguistic, $\Delta$=32.2pp) differ by only 4.2pp, with CP showing the lowest specification-layer rate and scopal the highest linguistic-layer rate. This overlap suggests that the layer boundary is not a sharp divide but a gradient, with CP and scopal occupying an intermediate region.

\subsection{C4: Model Variation}
\label{sec:model_variation}

Table~\ref{tab:model_layer} presents per-model layer gap analyses. Four of five models show a significant specification-over-linguistic gap (all $p<0.006$). The exception is Qwen3-8B, where the gap narrows to 4.3pp ($p$=0.281) and a partial inversion occurs: scopal ($\Delta$=32pp) exceeds CP ($\Delta$=20pp), reversing the within-layer ranking expected from the overall hierarchy.

\begin{table}[t]
\centering
\small
\caption{Per-model layer gap analysis. ``Spec Avg'' and ``Ling Avg'' are mean $\Delta$pp for specification- and linguistic-layer types.}
\label{tab:model_layer}
\begin{tabular}{lcccl}
\toprule
\textbf{Model} & \textbf{Spec} & \textbf{Ling} & \textbf{Gap} & \textbf{$p$} \\
\midrule
GPT-5.4 & 41.3 & 26.0 & +15.3 & $<$.001 \\
GPT-4.1 & 44.0 & 29.0 & +15.0 & $<$.001 \\
Claude Son.\ 4.6 & 39.0 & 25.0 & +14.0 & $<$.001 \\
Qwen3-32B & 36.7 & 24.7 & +12.0 & .006 \\
Qwen3-8B & 29.0 & 24.7 & +4.3 & .281 \\
\midrule
\textbf{Pooled} & \textbf{38.0} & \textbf{25.9} & \textbf{+12.1} & $<$\textbf{.001} \\
\bottomrule
\end{tabular}
\end{table}

The model$\times$type interaction is borderline significant ($\chi^2$=18.4, df=20, $p$=0.062, $V$=0.08), suggesting that the layer hierarchy is mostly universal but not perfectly clean. Logistic regression confirms a significant type main effect ($p<0.001$) but non-significant model main effect ($p$=0.34) and borderline interaction ($p$=0.058).

The Qwen3-8B deviation is attributable to its high baseline violation rate (19.2\% under unambiguous conditions, compared to 5.0\% for API models). This elevated baseline compresses relative differences between types, because the model fails frequently regardless of ambiguity type. Qwen3-32B (baseline 12.3\%) shows an intermediate pattern, maintaining a 12.0pp layer gap. This suggests that model capability moderates the clarity of the layer hierarchy: frontier models, with lower baselines, reveal the type structure more cleanly.

\subsection{API vs.\ Open-Source Models}

Table~\ref{tab:api_local} contrasts API and local model results. API models show lower baselines (5.0\% vs.\ 15.8\% unambiguous violation rate), larger binary effects (34.0pp vs.\ 30.0pp), and a clearer layer gap (14.8pp vs.\ 8.2pp). The API-only Cram\'{e}r's $V$=0.19 exceeds the pooled $V$=0.17, because open-source models' elevated baselines add noise that dilutes type-level differentiation.

\begin{table}[t]
\centering
\small
\caption{API (3 models) vs.\ local (2 models) comparison.}
\label{tab:api_local}
\begin{tabular}{lcc}
\toprule
\textbf{Metric} & \textbf{API} & \textbf{Local} \\
\midrule
Unamb.\ violation rate & 5.0\% & 15.8\% \\
Amb.\ violation rate & 39.0\% & 44.6\% \\
Binary $\Delta$pp & 34.0 & 30.0 \\
Layer gap (pp) & 14.8 & 8.2 \\
Cram\'{e}r's $V$ & 0.19 & 0.13 \\
\bottomrule
\end{tabular}
\end{table}

\subsection{Perturbation Control}
\label{sec:perturbation}

To distinguish ambiguity-specific effects from general prompt sensitivity, we construct 30 non-ambiguity perturbation pairs: synonym substitutions and clause reorderings that alter surface form without introducing ambiguity. Table~\ref{tab:perturbation} shows the results across 180 episodes (3 API models).

\begin{table}[t]
\centering
\small
\caption{Perturbation control results. ``Amb.\ ref.'' is the ambiguity effect from the main study for comparison.}
\label{tab:perturbation}
\begin{tabular}{lccc}
\toprule
\textbf{Condition} & \textbf{Viol.\ Rate} & \textbf{$\Delta$pp} & \textbf{95\% CI} \\
\midrule
Original & 5.6\% & -- & -- \\
Synonym subst. & 8.3\% & +2.8 & [$-$2.1, 7.7] \\
Clause reorder & 10.4\% & +4.8 & [$-$0.4, 10.0] \\
Combined & 9.4\% & +3.8 & [$-$0.9, 8.5] \\
\midrule
Amb.\ ref. & 39.0\% & +34.0 & [30.1, 37.9] \\
\bottomrule
\end{tabular}
\end{table}

The combined perturbation effect ($\Delta$=3.8pp) is an order of magnitude smaller than the ambiguity effect ($\Delta$=34.0pp), and its confidence interval contains zero. The ratio $\Delta_\text{perturb}/\Delta_\text{ambiguity}$=0.11. Non-ambiguity surface perturbations produce negligible violation rate changes, confirming that the observed 32pp effect is attributable to ambiguity rather than general prompt brittleness.

\section{Error Analysis}
\label{sec:error_analysis}

We manually categorize 300 violation traces (50 per type) into six failure modes. Table~\ref{tab:failure_modes} shows the distribution across types and layers.

\begin{table*}[t]
\centering
\small
\caption{Failure mode distribution (\% of violations per type). ``Assum.'' = assumption-based action, ``Scope'' = scope misapplication, ``Ref.'' = reference confusion, ``Arb.'' = arbitrary rule selection, ``Surf.'' = surface adoption, ``Cons.'' = conservative-but-wrong.}
\label{tab:failure_modes}
\begin{tabular}{lcccccc|cc}
\toprule
\textbf{Failure Mode} & \textbf{Auth.} & \textbf{Incomp.} & \textbf{CP} & \textbf{Scopal} & \textbf{Lex.} & \textbf{Coref.} & \textbf{Spec Avg} & \textbf{Ling Avg} \\
\midrule
Assumption-based & 71 & 58 & 42 & 28 & 32 & 18 & 62 & 35 \\
Scope misappl. & 2 & 4 & 18 & 38 & 22 & 8 & 8 & 32 \\
Reference confusion & 0 & 2 & 5 & 4 & 18 & 45 & 2 & 23 \\
Arbitrary rule sel. & 15 & 20 & 22 & 14 & 10 & 2 & 18 & 9 \\
Surface adoption & 4 & 8 & 8 & 12 & 14 & 0 & 7 & 10 \\
Conservative-wrong & 8 & 8 & 5 & 4 & 4 & 27 & 7 & 12 \\
\bottomrule
\end{tabular}
\end{table*}

\paragraph{Specification-layer failures.}
The dominant failure mode for specification-layer violations is assumption-based action (62\%): the agent fills an underspecified policy gap by adopting a ``reasonable'' assumption and acting on it, without any indication of uncertainty. Authorization scope violations are the most distinctive case, with 71\% involving unauthorized escalation, where the agent assumes the widest interpretation of ambiguous permission boundaries and executes actions beyond its intended authority. This pattern directly links ambiguity to safety violations through a causal pathway: vague permission language $\rightarrow$ assumption of broad authority $\rightarrow$ unauthorized action.

\paragraph{Linguistic-layer failures.}
Linguistic-layer violations show a more varied profile. Scope misapplication (32\%) is the predominant mode for scopal ambiguity (38\%), where the agent selects the wrong quantifier or negation scope and produces a ``good-enough'' but incorrect interpretation. Reference confusion (23\%) dominates coreferential violations (45\%), where the agent binds a relative clause to the nearest noun phrase without considering alternative attachments.

\paragraph{Predominant patterns, not clean dichotomy.}
The failure mode distribution shows predominant layer-level patterns, not a clean dichotomy. Assumption-based action, while dominant for specification-layer types (62\%), also accounts for 35\% of linguistic-layer violations, particularly for lexical ambiguity where the agent similarly lacks external disambiguation cues. Conversely, scope misapplication occurs in 18\% of CP violations. The layer-level failure mode match rate is 62\%, with 38\% cross-layer overlap.

\paragraph{Coreferential avoidance.}
Among coreferential violations, 28\% exhibit avoidance behavior: the agent shows partial awareness of the ambiguity (e.g., briefly hesitating or requesting clarification) but ultimately commits to an incorrect entity binding. This contrasts with the 45\% showing pure reference confusion with no apparent awareness. The avoidance pattern provides suggestive evidence for the disambiguation channel hypothesis, as agents occasionally access contextual cues for referential ambiguity, but the evidence is insufficient for a standalone claim.

\paragraph{Tool-based partial disambiguation.}
Across all types, 23\% of violation traces show evidence of partial disambiguation through tool responses. This occurs predominantly in linguistic-layer episodes (34\% vs.\ 8\% for specification-layer), supporting the hypothesis that linguistic ambiguity is more amenable to contextual resolution. However, partial disambiguation does not guarantee correct resolution: approximately half of these attempts still result in violations.

\subsection{Inter-Model Nocuity}
\label{sec:nocuity}

Following Berry and Kamsties' (\citeyear{berry2003ambiguity}) definition, we examine whether different models produce \emph{divergent} interpretations of the same ambiguous clause or converge on the same (wrong) interpretation. Fleiss' $\kappa$ across five models on ambiguous clauses is 0.44 (moderate agreement). Of all ambiguous clauses, 38\% elicit unanimous (but incorrect) interpretations across all models, 42\% produce genuinely divergent interpretations ($\geq$2 distinct wrong readings), and 20\% show mixed patterns. The 42\% divergent rate constitutes true nocuity in Berry and Kamsties' sense, while the 38\% convergent-but-wrong rate represents a distinct phenomenon: systematic bias, where all agents adopt the same surface interpretation despite it being incorrect.

Linguistic-layer clauses show higher inter-model agreement ($\kappa$=0.51) than specification-layer clauses ($\kappa$=0.39), suggesting that linguistic ambiguity more often triggers a single dominant parse across models, while specification underspecification admits more varied ``reasonable'' assumptions.

\section{Judge Calibration}
\label{sec:judge_calibration}

We validate the LLM judge against 100 episodes annotated by two human experts (inter-annotator $\kappa$=0.78). Table~\ref{tab:judge_calibration} reports per-type agreement.

\begin{table}[t]
\centering
\small
\caption{LLM judge agreement with human annotators.}
\label{tab:judge_calibration}
\begin{tabular}{lcccc}
\toprule
\textbf{Type} & \textbf{F1} & \textbf{Prec.} & \textbf{Rec.} & \textbf{Hum.\ $\kappa$} \\
\midrule
Auth.\ scope & 0.83 & 0.86 & 0.80 & 0.82 \\
Incompleteness & 0.82 & 0.84 & 0.80 & 0.80 \\
CP & 0.78 & 0.80 & 0.76 & 0.78 \\
Scopal & 0.81 & 0.83 & 0.79 & 0.79 \\
Lexical & 0.72 & 0.75 & 0.69 & 0.72 \\
Coreferential & 0.71 & 0.74 & 0.68 & 0.70 \\
\midrule
\textbf{Overall} & \textbf{0.80} & \textbf{0.83} & \textbf{0.77} & \textbf{0.78} \\
\bottomrule
\end{tabular}
\end{table}

Overall F1 is 0.80 ($\kappa$=0.73). Two types fall below the 0.75 F1 threshold: lexical (0.72) and coreferential (0.71). For lexical ambiguity, the violation boundary is inherently subjective, as polysemy produces interpretations that may be ``non-intended but reasonable'' rather than clearly wrong. For coreferential ambiguity, avoidance behavior blurs the violation boundary when agents adopt cautious strategies that still produce incorrect outcomes. Performance also varies by severity: clear violations (F1=0.87) and non-violations (F1=0.85) are well-identified, while borderline cases (F1=0.68) challenge both the judge and human annotators. All main findings should be interpreted in light of this measurement imprecision, though the binary effect ($\Delta$=32pp) far exceeds the judge error margin.

\section{Practical Application: Ambiguity Detection}
\label{sec:classifier}

As a proof-of-concept, we fine-tune a DeBERTa-v3-large binary ambiguity classifier on the 300 clause pairs (240 train / 60 test). Table~\ref{tab:classifier} reports per-type results.

\begin{table}[t]
\centering
\small
\caption{DeBERTa-v3-large binary ambiguity classifier.}
\label{tab:classifier}
\begin{tabular}{lccc}
\toprule
\textbf{Type} & \textbf{F1} & \textbf{Prec.} & \textbf{Rec.} \\
\midrule
Auth.\ scope & 0.86 & 0.88 & 0.84 \\
Incompleteness & 0.83 & 0.85 & 0.81 \\
CP & 0.79 & 0.82 & 0.76 \\
Scopal & 0.76 & 0.78 & 0.74 \\
Lexical & 0.73 & 0.76 & 0.70 \\
Coreferential & 0.68 & 0.72 & 0.64 \\
\midrule
\textbf{Overall} & \textbf{0.78} & \textbf{0.81} & \textbf{0.75} \\
Cross-type transfer & 0.65 & -- & -- \\
\bottomrule
\end{tabular}
\end{table}

The classifier achieves F1=0.78 overall, with a notable parallel between detection ease and danger: the most detectable types (authorization scope F1=0.86, incompleteness F1=0.83) are also the most dangerous. This suggests that even a simple binary classifier can prioritize the highest-risk ambiguity types for human review. Cross-type transfer F1=0.65 (training on five types, testing on the held-out type) indicates that type-agnostic detection remains challenging and that type-specific training data improves performance substantially.

\section{Discussion}

\paragraph{Differential nocuity as a framework.}
Our results establish that all six ambiguity types are nocuous for LLM agents, confirming the hypothesis that agents lack the shared organizational context that renders most ambiguity innocuous for human readers \citep{berry2003ambiguity}. However, nocuity is not uniform. Specification-layer ambiguities produce substantially higher violation rates than linguistic-layer ambiguities, with a layer gap of 12.1pp (CI: [6.8, 17.4]). We propose the term \emph{differential nocuity} to describe this phenomenon: all policy ambiguity types are nocuous, but the severity of safety consequences depends systematically on the ambiguity's location in the specification-to-language hierarchy.

\paragraph{Disambiguation channel hypothesis.}
The layer hierarchy reflects the availability of external disambiguation resources. When an agent encounters a linguistically ambiguous clause (e.g., scopal or coreferential), tool responses, user messages, and world knowledge sometimes provide cues sufficient for partial resolution (observed in 34\% of linguistic-layer violation traces). Specification-layer ambiguity (e.g., undefined permission boundaries, missing exception conditions) has no such external resolution channel: the information simply does not exist in the agent's environment, forcing assumption-based action. This asymmetry in disambiguation channels explains why specification-layer violations are more frequent, not because the ambiguity is more confusing, but because agents have no path to resolving it.

\paragraph{Model capability as moderator.}
The layer hierarchy is robust across frontier models (gap 14.0 to 15.3pp) but attenuated in smaller open-source models. Qwen3-8B shows a gap of only 4.3pp with a partial layer inversion (scopal $\Delta$=32pp $>$ CP $\Delta$=20pp). This attenuation is attributable to the model's elevated baseline violation rate (19.2\% under unambiguous conditions), which compresses relative differences between types. When a model violates policies frequently regardless of ambiguity, the additional risk from specification-layer ambiguity is proportionally smaller. The borderline model$\times$type interaction ($p$=0.062) confirms that the layer effect is mostly universal but not perfectly clean, with smaller models deviating from the pattern.

\paragraph{Convergence zone: CP and scopal.}
CP ($\Delta$=28.0pp, specification) and scopal ($\Delta$=32.2pp, linguistic) differ by only 4.2pp, forming a convergence zone at the layer boundary. CP shows the lowest specification-layer rate because conditional presuppositions sometimes have detectable surface cues (explicit ``if'' clauses, numbered conditions). Scopal ambiguity shows the highest linguistic-layer rate because quantifier scope rarely has external resolution cues. This overlap does not invalidate the layer distinction but highlights that the boundary is a gradient rather than a sharp divide.

\paragraph{Practical implications.}
The differential nocuity finding motivates a tiered auditing strategy for agent policies. Authorization scope and incompleteness ambiguities, which produce the highest violation rates and are the most detectable by automated classifiers, should receive priority review. Our proof-of-concept classifier (F1=0.78) demonstrates that automated pre-screening is feasible, though type-specific training data is needed for acceptable performance on all types.

\subsection{Limitations}

\paragraph{Judge calibration.}
Lexical (F1=0.72) and coreferential (F1=0.71) judge F1 scores fall below the 0.75 threshold, indicating that violation boundaries for these types are inherently fuzzy. Borderline-severity episodes are identified at only F1=0.68. While the main binary effect ($\Delta$=32pp) far exceeds the judge error margin, fine-grained type rankings within a layer may be influenced by differential judge accuracy.

\paragraph{Synthetic clause construction.}
Our clauses are somewhat shorter and less lexically diverse than naturally occurring policy ambiguities (Table~\ref{tab:realism}). The matched-pair design necessitates controlled construction, which inherently constrains surface properties. However, the structural property most relevant to ambiguity processing (parse depth) shows no distributional difference, and the 32pp violation effect is unlikely to be driven by surface-level features.

\paragraph{Domain coverage.}
$\tau^2$-bench covers retail and airline domains. Whether the layer hierarchy generalizes to other policy domains (healthcare, legal, financial) remains an open question.

\paragraph{Model$\times$type interaction.}
The borderline interaction ($p$=0.062) means we cannot cleanly claim that the layer hierarchy is universal across all model scales. The Qwen3-8B deviation suggests that weaker models may not preserve the same type-level structure observed in frontier models.

\paragraph{Error analysis overlap.}
Failure modes show predominant layer-level patterns but substantial cross-layer overlap (38\%). The error taxonomy should be interpreted as capturing tendencies rather than deterministic type-to-failure mappings.

\section{Conclusion}

We present the first causal, type-granular study of how policy ambiguity affects safety violations in LLM agents. Across 3,000 episodes, five models, and six ambiguity types, ambiguous policy clauses increase violation rates by 31.9 percentage points. This effect is attributable to ambiguity rather than general prompt sensitivity, as confirmed by perturbation controls.

Not all ambiguity is equally dangerous. Specification-layer ambiguities (authorization scope, incompleteness, conditional presupposition) produce substantially higher violation rates than linguistic-layer ambiguities (scopal, lexical, coreferential), with a pooled gap of 12.1pp. This differential nocuity reflects the availability of external disambiguation channels: linguistic ambiguity can sometimes be resolved through contextual cues, while specification underspecification forces agents into assumption-based action.

These findings motivate type-aware policy auditing that prioritizes specification-layer ambiguities in safety-critical deployments. We release all 300 clause pairs, evaluation scripts, and analysis code to support future work on ambiguity-aware agent safety.

\bibliography{references}
\bibliographystyle{acl_natbib}

\appendix

\section{Mini-Pilot Results}
\label{app:mini_pilot}

Prior to the full study, we conducted a mini-pilot of the two new ambiguity types (authorization scope and coreferential) on GPT-5.4: 5 clauses per type $\times$ 2 conditions = 20 episodes. Authorization scope showed a 60\% ambiguous violation rate (10\% unambiguous), confirming strong signal. Coreferential showed 30\% ambiguous (10\% unambiguous), indicating weaker but sufficient signal. Two coreferential clauses were flagged for revision: one with overly obvious RC attachment (agent resolved with zero hesitation) and one where the unambiguous wording introduced new ambiguity. Both were revised before the full study.

\section{Detailed Per-Model Results}
\label{app:per_model}

Table~\ref{tab:api_per_type} presents the full per-model, per-type ambiguous violation rate matrix for API models.

\begin{table*}[t]
\centering
\small
\caption{Per-model ambiguous violation rate (\%) for API models.}
\label{tab:api_per_type}
\begin{tabular}{lcccc}
\toprule
\textbf{Type} & \textbf{GPT-5.4} & \textbf{GPT-4.1} & \textbf{Claude Son.\ 4.6} & \textbf{3-API Avg} \\
\midrule
Auth.\ scope & 52 & 56 & 50 & 52.7 \\
Incompleteness & 48 & 53 & 46 & 49.0 \\
CP & 33 & 36 & 30 & 33.0 \\
Scopal & 38 & 42 & 36 & 38.7 \\
Lexical & 35 & 40 & 33 & 36.0 \\
Coreferential & 24 & 28 & 22 & 24.7 \\
\midrule
\textbf{Overall} & \textbf{38.3} & \textbf{42.5} & \textbf{36.2} & \textbf{39.0} \\
\bottomrule
\end{tabular}
\end{table*}

Table~\ref{tab:local_per_type} presents results for local (open-source) models.

\begin{table*}[t]
\centering
\small
\caption{Per-model ambiguous violation rate (\%) and $\Delta$pp for local models.}
\label{tab:local_per_type}
\begin{tabular}{lccc|cc}
\toprule
\textbf{Type} & \textbf{Qwen3-32B Amb} & \textbf{Qwen3-8B Amb} & \textbf{Local Avg} & \textbf{Qwen3-32B $\Delta$} & \textbf{Qwen3-8B $\Delta$} \\
\midrule
Auth.\ scope & 54 & 55 & 54.5 & +42 & +35 \\
Incompleteness & 50 & 50 & 50.0 & +40 & +32 \\
CP & 38 & 38 & 38.0 & +28 & +20 \\
Scopal & 42 & 50 & 46.0 & +28 & +32 \\
Lexical & 40 & 46 & 43.0 & +27 & +26 \\
Coreferential & 34 & 38 & 36.0 & +19 & +16 \\
\midrule
\textbf{Overall} & \textbf{43.0} & \textbf{48.5} & \textbf{44.6} & \textbf{+30.7} & \textbf{+29.3} \\
\bottomrule
\end{tabular}
\end{table*}

\section{Qwen3-8B Layer Inversion Detail}
\label{app:qwen_inversion}

Qwen3-8B shows a partial layer inversion where scopal (linguistic, $\Delta$=32pp) exceeds CP (specification, $\Delta$=20pp) by 12pp. Table~\ref{tab:qwen_inversion} details the cross-layer comparisons.

\begin{table}[t]
\centering
\small
\caption{Qwen3-8B cross-layer type comparisons.}
\label{tab:qwen_inversion}
\begin{tabular}{llcl}
\toprule
\textbf{Spec Type ($\Delta$)} & \textbf{Ling Type ($\Delta$)} & \textbf{Diff} & \textbf{Ling$>$Spec?} \\
\midrule
CP (20) & Scopal (32) & +12 & Yes \\
CP (20) & Lexical (26) & +6 & Yes \\
Incomp.\ (32) & Scopal (32) & 0 & Tie \\
Auth.\ (35) & Scopal (32) & $-$3 & No \\
\bottomrule
\end{tabular}
\end{table}

\end{document}
